\section{Design}

This chapter describes the design decisions and methods I will adopt and have adopted so far in this project. Building on the limitations identified in the literature review, the design priorities, reproducibility, and clinical relevance over maximising headline performance metrics. The aim is not to propose an autonomous diagnostic system, but to design and evaluate a deep learning-based decision support tool that could realistically assist clinicians in breast cancer screening workflows. Whilst there will be a high level of automation, the idea will be that the system will highlight mammograms of concern for further inspection. This still, in principle, speeds up the clinician’s workload, not having to analyse every single mammography image as it will automatically sort those with or without cancer at a high certainty, but it will also allow those mammograms within a grey area of uncertainty to be evaluated further.

The overall system follows a supervised deep learning pipeline for binary classification of mammography images into cancerous and non-cancerous classes. Design choices regarding dataset selection, preprocessing, model architecture, training strategy, and evaluation metrics are explicitly informed by findings from prior work, particularly with respect to data leakage, dataset imbalance, and the clinical implications of false-negative predictions with extra care taken to address data imbalance within the available datasets.

The Curated Breast Imaging Subset of the Digital Database for Screening Mammography (CBIS-DDSM) was selected as the primary dataset for this project. CBIS-DDSM is a widely used benchmark dataset that provides curated mammography images with standardised labels and associated metadata, including patient identifiers (Lee et al., 2017). The availability of patient IDs is a critical factor, as it enables patient-level dataset splitting rather than image-level splitting. As highlighted by Wang et al. (2024), many mammography studies report inflated performance due to information leakage when images from the same patient appear in both training and test sets. By enforcing patient-level separation, this project aims to produce a more realistic estimate of generalisation performance.

Before training the model, the mammography images are preprocessed to ensure they are consistent and suitable for input into a convolutional neural network. Images are resized to a fixed resolution and normalised so that pixel intensity values are comparable across the dataset, which helps the network learn more stable feature representations during training. Given that mammograms are high-resolution grayscale images, excessive downsampling is avoided, as this could remove fine textural details that may be important for identifying early-stage abnormalities.

During training, data augmentation is applied using simple transformations such as horizontal flipping, small rotations, and mild contrast adjustments. These techniques are used to expose the model to a wider range of plausible image variations and reduce overfitting, without changing the underlying medical content of the images. This approach is commonly adopted in medical imaging studies where labelled data is limited and has been shown to improve generalisation when applied carefully (Ragab et al., 2022).

The main part of the system is a CNN, as they are most effective in image-based diagnosis. CNNs are well suited to mammography analysis because they learn hierarchical feature representations directly from pixel data, allowing early layers to capture low-level features such as edges and textures, while deeper layers encode higher level semantic structures associated with lesions or abnormal tissue patterns. This hierarchical learning capability has been shown to outperform traditional computer aided detection systems based on handcrafted features, as discussed in the earlier literature review (Wang et al., 2024).

The baseline architecture follows a standard CNN design consisting of stacked convolutional layers, nonlinear activation functions, and pooling layers. Convolutional layers apply learnable filters across the image to extract spatial features, while rectified linear unit (ReLU) activations introduce nonlinearity, enabling the network to model complex relationships within the data. Max-pooling layers are used intermittently to reduce spatial dimensionality and improve translation invariance, which is particularly important in mammography where lesions may appear at varying locations within the breast tissue, although it is worth noting, that the placement of some lesions have been researched into how likely a cancer is to spread (Yang et al); however, I have decided not to factor this into my design as it adds an extra layer of complication and is not an exact measure of the likelihood a cancer will metastise.

To address the challenge of limited labelled data, my project will incorporate transfer learning using a pre-trained convolutional network. The literature shows that transfer learning is becoming widely adopted in medical imaging to mitigate data scarcity by leveraging representations learned from large-scale natural image datasets such as ImageNet (Ragab et al., 2022). In my project, a pre-trained ResNet-based architecture is used as the feature extractor, with the final classification layers retrained on the mammography dataset. Although natural images differ significantly from medical images, prior studies suggest that early convolutional layers capture generic visual features that can be effectively repurposed for medical tasks when fine-tuned appropriately (Ragab et al., 2022).

The architecture is adapted by replacing the original classification head with a custom fully connected layer tailored to binary classification. Global average pooling is employed prior to the final dense layer to reduce the number of trainable parameters and lower the risk of overfitting. Again, during the literature review, overfitting was an issue that arose on a few occasions, so it is best to take steps to mitigate this. This design choice is particularly important given the imbalance between positive cancer and negative cancer cases in screening datasets. Dropout regularisation is applied during training to further improve generalisation by discouraging reliance on any single feature pathway.

Model training is performed using a binary cross-entropy loss function, which is appropriate for binary classification tasks and aligns naturally with probabilistic outputs from the final sigmoid layer. The Adam optimiser is used to update model weights during training, as it adapts learning rates automatically for individual parameters, helping to stabilise training and reduce sensitivity to initial hyperparameter choices.

Training is carried out over multiple epochs, with early stopping applied based on validation performance to reduce the risk of overfitting. Key hyperparameters such as learning rate, batch size, and dropout probability are selected empirically through initial experimentation rather than exhaustive optimisation. I have opted for this approach to take a more exploratory route in order to prioritise understanding model behaviour over fine-tuning for marginal performance gains.

A key design priority is the handling of class imbalance, which is inherent in mammography datasets where malignant cases are significantly rarer than benign or normal cases. Relying solely on accuracy in such scenarios can be misleading, as a model may achieve high accuracy by predominantly predicting the majority class. To mitigate this, class weighting is applied during training to penalise misclassification of cancer-positive cases more heavily. This aligns with the clinical objective of minimising false negatives, as missed diagnoses can have severe consequences, an ethical consideration that must be accounted for in a project that has medical ramifications for patients.

Evaluation strategy is a central component of the system design and is directly informed by shortcomings identified in prior work. While many studies report area under the receiver operating characteristic curve (AUC) as the primary performance metric, AUC alone does not fully capture clinical utility, particularly in imbalanced datasets (Saito and Rehmsmeier, 2015). In this project, AUC is reported to allow comparison with existing literature, but it is complemented by additional metrics that better reflect real-world screening priorities.

These metrics include sensitivity (recall), specificity, precision, and F1-score. Sensitivity is particularly emphasised, as it measures the proportion of actual cancer cases correctly identified by the model. In screening contexts, high sensitivity is often prioritised over specificity to reduce the risk of missed cancers. Precision and F1-score provide insight into the trade-off between false positives and false negatives, offering a more balanced assessment than accuracy alone (Saito and Rehmsmeier, 2015). From my perspective, I feel that it would be better to be overcautious than the reverse; a false positive, whilst inconvenient for a number of reasons to a fictitious patient, is better than a false negative resulting in a missed opportunity to diagnose and treat cancer as soon as possible. As per my previous research, it is clear that the sooner a positive diagnosis is made, the better chance of a positive outcome for that patient.

This ties in to how my project will aim to look beyond fixed threshold metrics, the project will explore threshold-based decisions designed to support a human-in-the-loop deployment model. Rather than treating the CNN as an autonomous classifier, predicted probabilities are interpreted as confidence scores. High-confidence positive predictions may be flagged for urgent clinical review, while low-confidence cases may be deferred to human experts. This approach aligns with recent recommendations for safe medical AI deployment, which emphasise decision support rather than full automation (Wang et al., 2024).

Patient level evaluation is strictly enforced throughout model validation and testing. Data splits will be performed such that the aim will be that no patient appears in more than one subset, ensuring that performance metrics reflect true generalisation to unseen individuals. This design choice directly addresses data leakage concerns raised in the literature review and strengthens the credibility of reported results.

Reproducibility is also considered an essential design requirement. All experiments are conducted within a version-controlled codebase hosted on GitHub, with incremental commits documenting model changes, parameter adjustments, and evaluation outcomes. Random seeds are fixed where possible to reduce nondeterminism, and experiment configurations are logged to enable consistent replication of results, this was an error I made on my previous Machine Learning and Neural Network project, something I am aiming to avoid on this project. This emphasis on transparency and traceability reflects best practice in modern machine learning research and aligns with assessment requirements for the project. 

In summary, the system design adopts a cautious and clinically informed approach to deep learning-based breast cancer detection. By combining convolutional neural networks with transfer learning, patient level validation, and clinically relevant evaluation metrics, the project seeks to address key limitations identified in prior work. Rather than optimising solely for headline performance, the design prioritises robustness, interpretability, and alignment with real-world screening needs, providing a foundation for meaningful evaluation and future refinement.
